%% kapa: exam=1 round=19 year=2008 period=2 form=A subject=economics subject_ko=경제원론 seq=28 number=28
%% source: Q-Net 감정평가사 기출문제 (cst00309) artlSeq=1001962 file=extracted/1008/1001962/2008년 제19회 감정평가사 1차시험 기출문제/2교시.hwp
%% conversion: kordoc 4.24.51 -> tools/gen1.py (auto); review=proofread-sonnet; flags=none
\begin{question}{28}
케인즈 단순모형에 기초한 총수요․총공급 모형에서 경제는 완전고용상태에 있다고 하자. 정부지출을 증가시키면서 총수요 증가로 인한 물가 상승을 유발하지 않으려면 조세를 어떻게 해야 하는가? (단, 단기총공급곡선은 우상향한다고 가정한다.)
\begin{choices}
  \item 줄여야 한다.
  \item 정부지출보다 적게 늘려도 된다.
  \item 정부지출만큼 늘리면 된다.
  \item 정부지출보다 많이 늘려야 한다.
  \item 어떤 조세정책을 써도 물가 상승이 유발된다.
\end{choices}
\end{question}
\answer{4}
